\chapter{Evaluation}
\label{ch:evaluation}
\noindent
This chapter evaluates the results reported in Chapter~\ref{ch:results}
against the requirements defined in Chapter~3. The results chapter shows what
the implemented architecture produced. This chapter considers whether those
outputs satisfy the requirements and whether they satisfy the auditability
definition used throughout the thesis.

The evaluation uses the definition from Chapter~2, where auditability is framed
as the designed capability of a system to allow its processes, controls, and
outcomes to be observed and verified by stakeholders. In that definition,
auditability is not a single log file or a general claim that monitoring
exists. It is made up of normativity, accountability, referentiality
infrastructure, reliability, and verifiability
\cite{weigand_auditability_2019}. The evaluation therefore asks whether the
implemented fields and artifacts make those dimensions visible in the system.

\section{Evaluation Approach}
\noindent
The assessment follows the R1--R5 requirement matrix from Chapter~3 and the
control matrix in the proposed-solution repository. Each requirement is
evaluated in three steps. First, the requirement is mapped back to the relevant
auditability dimensions. Second, the concrete evidence fields are identified.
Third, the observed result is judged against the requirement.

This approach deliberately treats the evidence artifacts as the primary object
of evaluation. A pipeline run is not considered auditable only because it
completed successfully. It is considered auditable if the produced dataset can
be linked to the governing change, release, infrastructure reference, runtime
code, quality decision, lineage event, and metadata record without manual log
collation.

\section{R1 Governance Accountability}
\noindent
R1 requires production-relevant changes to be linked to documented authority.
In Chapter~3 this requirement was tied to normativity and accountability.
Normativity is satisfied when the system records which governance rule or
change process applies to an action. Accountability is satisfied when the run
can be linked back to the actor, change, or approved release that authorised
the action.

\begin{table}[H]
  \centering
  \caption{R1 governance accountability evidence fields.}
  \label{tab:evaluation-r1-fields}
  \renewcommand{\arraystretch}{1.15}
  \begin{tabularx}{\textwidth}{@{}p{0.33\textwidth}X@{}}
    \toprule
    \textbf{Field or artifact} & \textbf{Evaluation role} \\
    \midrule
    \texttt{change\_ref} & Links the run to the governing change record. \\
    \texttt{release\_manifest\_ref} & Links the run to the approved release definition. \\
    \texttt{terraform\_state\_ref} & Links the run to the infrastructure state reference. \\
    \texttt{reference\_checks.change\_ref} & Records whether the change record resolved at runtime. \\
    \texttt{reference\_checks.release\_manifest\_ref} & Records whether the release manifest resolved at runtime. \\
    \texttt{reference\_checks.terraform\_state\_ref} & Records whether the infrastructure reference resolved at runtime. \\
    \texttt{source\_control.commit} & Links the release and run evidence to the source revision. \\
    Change record \texttt{change\_id} & Names the governed change, for example \texttt{EDGAR-1}. \\
    Release manifest \texttt{release\_id} & Names the governed release, for example \texttt{dev-edgar-v1}. \\
    \bottomrule
  \end{tabularx}
\end{table}

The requirement is demonstrated because the implemented run is not recorded as
an isolated Airflow execution. The run manifest contains the change, release,
and Terraform references, and the reference-check block records whether each
reference existed when the run evidence was produced. This means the run can be
read as an authorised operational event rather than only as a scheduled task.

The link to the auditability definition is direct. The \texttt{change\_ref} and
\texttt{release\_manifest\_ref} fields provide normativity because they record
the governing process that applied to the run. The source-control and release
fields provide accountability because they bind the run to a concrete change
record and code revision. The reference checks are important because they
prevent the evidence chain from merely naming authority; they record that the
named authority resolved.

\section{R2 Information Security Assurance}
\noindent
R2 requires deployable artifacts and configuration to be verified for integrity
and security before they are used. In Chapter~3 this was tied to reliability
and verifiability. Reliability is concerned with whether the evidence faithfully
represents the artifact that ran. Verifiability is concerned with whether an
independent reviewer can test the claim.

\begin{table}[H]
  \centering
  \caption{R2 information security assurance evidence fields.}
  \label{tab:evaluation-r2-fields}
  \renewcommand{\arraystretch}{1.15}
  \begin{tabularx}{\textwidth}{@{}p{0.33\textwidth}X@{}}
    \toprule
    \textbf{Field or artifact} & \textbf{Evaluation role} \\
    \midrule
    \texttt{code\_bundle.release\_manifest\_sha256} & Records the digest of the release manifest. \\
    \texttt{code\_bundle.artifacts.path} & Names each DAG or job artifact included in the bundle. \\
    \texttt{code\_bundle.artifacts.sha256} & Records a SHA-256 hash for each code artifact. \\
    \texttt{source\_control.repository} & Records the source repository used for the release. \\
    \texttt{source\_control.commit} & Records the source revision associated with the run. \\
    \texttt{attestation.provider} & Records the attestation provider expected by the evidence model. \\
    \texttt{attestation.subject\_digest} & Carries the subject digest used to bind provenance to the release artifact. \\
    GitHub Actions controls & Provide the CI/CD surface for linting, tests, policy checks, artifact upload, and attestation. \\
    \bottomrule
  \end{tabularx}
\end{table}

The requirement is demonstrated by the way the implementation records artifact
identity. The code bundle includes hashes for the DAG and the ingest,
transform, quality, and promotion jobs. These hashes make the run evidence more
reliable because an auditor can compare the recorded digest with the file that
is claimed to have run. The source-control fields then bind those artifacts to
the repository and commit.

The release provenance fields extend the same chain to the release context. The
evidence model records the provider, workflow, run, URL, and subject digest
fields needed to connect the release bundle to GitHub-native provenance. In
the implemented results, this requirement is evaluated through the demonstrated
source revision, run reference, release reference, and SHA-256 artifact hashes.

This satisfies the auditability definition because the security claim is
independently testable. An evaluator does not need to trust a written statement
that the correct files were used. The evaluator can inspect
\texttt{code\_bundle.artifacts.sha256}, compare it with the artifact contents,
and then follow \texttt{source\_control.commit} back to the release source.

\section{R3 Operational Risk Control}
\noindent
R3 requires the workflow to support reproducibility, recovery, and visibility
over operational state. In Chapter~3 this requirement was tied to
referentiality and reliability. Referentiality is satisfied when the run record
connects the output to its source, intermediate, and runtime records.
Reliability is satisfied when the recorded state is complete enough to support
replay and diagnosis.

\begin{table}[H]
  \centering
  \caption{R3 operational risk control evidence fields.}
  \label{tab:evaluation-r3-fields}
  \renewcommand{\arraystretch}{1.15}
  \begin{tabularx}{\textwidth}{@{}p{0.33\textwidth}X@{}}
    \toprule
    \textbf{Field or artifact} & \textbf{Evaluation role} \\
    \midrule
    \texttt{run\_id} & Identifies the specific governed execution. \\
    \texttt{pipeline\_id} & Identifies the pipeline definition that ran. \\
    \texttt{dataset\_date} & Identifies the dataset partition or business date. \\
    \texttt{source\_ref} & Records the source input used by the run. \\
    \texttt{raw\_outputs} & Records the raw payload and normalised raw output references. \\
    \texttt{curated\_outputs.curated\_uri} & Records the curated dataset produced by transformation. \\
    \texttt{quality\_result.status} & Records whether the quality gate passed. \\
    \texttt{gold\_write\_result} & Records the promoted gold table write result. \\
    \texttt{dataset\_version} & Records the stable output version used for reconstruction. \\
    \texttt{status} & Records the final run outcome. \\
    \texttt{completed\_at} & Records when the run evidence was finalised. \\
    \bottomrule
  \end{tabularx}
\end{table}

The requirement is demonstrated because the run manifest records enough state
to reconstruct how the output was created. It contains source, raw, curated,
quality, promotion, and dataset-version fields in a single evidence object.
This reduces operational risk because recovery no longer depends on reading
unrelated logs and guessing which task output should be trusted.

The manifest validator strengthens the result. For a successful run, the
validator requires a \texttt{dataset\_version}, a \texttt{gold\_write\_result},
and metadata references to the OpenMetadata run record, OpenMetadata dataset
record, OpenLineage event, and audit-chain path. It also requires reference
checks for the release, Terraform, and change references. A run is therefore
not merely successful because the final task completed. It is successful
because the evidence needed for reconstruction is present.

This maps to referentiality because each stage of the data path points to the
next stage. It maps to reliability because the system records the final status,
quality result, row-count evidence, and completion time together with the
output version. The result is a recovery-oriented run record rather than a
best-effort operational log.

\section{R4 Continuous Control Monitoring}
\noindent
R4 requires controls to run continuously and to produce evidence of their
outcomes. In Chapter~3 this requirement was tied to normativity and
accountability. Normativity is present when controls are embedded into the
workflow rather than applied after the fact. Accountability is present when the
result of each control decision is retained.

\begin{table}[H]
  \centering
  \caption{R4 continuous control monitoring evidence fields.}
  \label{tab:evaluation-r4-fields}
  \renewcommand{\arraystretch}{1.15}
  \begin{tabularx}{\textwidth}{@{}p{0.33\textwidth}X@{}}
    \toprule
    \textbf{Field or artifact} & \textbf{Evaluation role} \\
    \midrule
    \texttt{quality\_result.checks.curated\_output\_exists} & Records whether the curated output existed. \\
    \texttt{quality\_result.checks.schema\_matches\_contract} & Records whether the curated output matched the contract. \\
    \texttt{quality\_result.checks.row\_count\_gt\_zero} & Records whether the curated output contained records. \\
    \texttt{quality\_result.checks.required\_ids\_non\_null} & Records whether required identifiers were populated. \\
    \texttt{curated\_outputs.contract\_path} & Links the quality decision to the data contract. \\
    \texttt{reference\_checks} & Records whether governance references resolved before the run evidence was accepted. \\
    \texttt{metadata\_refs.openlineage\_event\_path} & Records where the lineage event evidence was written. \\
    \texttt{metadata\_refs.openmetadata\_run\_path} & Records where the run metadata evidence was written. \\
    \bottomrule
  \end{tabularx}
\end{table}

The requirement is demonstrated because the implementation converts controls
into runtime evidence. The quality result is not a general pass/fail statement;
it records named checks. The contract path also remains in the manifest, which
means the schema decision can be traced back to the exact contract used for the
check.

This matters for auditability because continuous monitoring must be observable.
If a quality or governance check only appears in a log line, the control is
difficult to test independently. In this implementation, the check result is a
field in the manifest and therefore becomes part of the same evidence chain as
the dataset version. The same pattern is used for governance references:
\texttt{reference\_checks} records whether the release manifest, Terraform
state reference, and change record resolved.

R4 is therefore satisfied within the implemented scope. The system demonstrates
that control execution and control outcome can be captured as structured
runtime evidence. The remaining production concern is scale and policy
coverage: more policies would be needed for a complete enterprise deployment,
but the architecture pattern is demonstrated.

\section{R5 End-to-End Evidence Chain}
\noindent
R5 is the strongest auditability requirement because it asks whether a promoted
dataset can be reconstructed from evidence without manual log collation. In
Chapter~3 this requirement was tied to referentiality and verifiability. It is
also the requirement that most directly tests the thesis claim.

\begin{table}[H]
  \centering
  \caption{R5 end-to-end evidence chain fields.}
  \label{tab:evaluation-r5-fields}
  \renewcommand{\arraystretch}{1.15}
  \begin{tabularx}{\textwidth}{@{}p{0.33\textwidth}X@{}}
    \toprule
    \textbf{Field or artifact} & \textbf{Evaluation role} \\
    \midrule
    \texttt{dataset\_version} & Starting point for reconstruction from the gold dataset. \\
    Audit-chain \texttt{manifest\_ref} & Resolves the dataset version to the run manifest. \\
    Audit-chain \texttt{release\_manifest\_ref} & Resolves the dataset version to the release manifest. \\
    Audit-chain \texttt{terraform\_state\_ref} & Resolves the dataset version to the infrastructure state reference. \\
    Audit-chain \texttt{change\_ref} & Resolves the dataset version to the governing change. \\
    Audit-chain \texttt{code\_bundle} & Carries the runtime code hashes into the reconstruction path. \\
    \texttt{metadata\_refs.audit\_chain\_path} & Records where the audit-chain artifact was written. \\
    OpenLineage \texttt{inputs} & Records source, raw, curated, and manifest lineage inputs. \\
    OpenLineage \texttt{outputs} & Records the gold-table lineage output and dataset version facet. \\
    OpenMetadata dataset \texttt{manifest\_ref} & Links the catalogue dataset record back to the run manifest. \\
    \bottomrule
  \end{tabularx}
\end{table}

The implemented result satisfies R5 because the dataset version
\texttt{9b5fd26a2005c02b} can be used as the first lookup key. From there, the
audit-chain artifact points to the run manifest. The run manifest then points
to the release manifest, Terraform state reference, change record, code bundle,
quality result, OpenLineage event, and OpenMetadata records.

This is the clearest link back to the auditability definition. Referentiality
is demonstrated because the evidence chain preserves the points of recording
for source data, raw output, curated output, manifest storage, metadata
publication, and gold promotion. Verifiability is demonstrated because those
links are materialised as fields and paths rather than implied by naming
convention. An evaluator can follow the fields directly:
\texttt{dataset\_version} to audit-chain file, audit-chain
\texttt{manifest\_ref} to run manifest, run manifest to governance references,
and metadata references to lineage and catalogue records.

R5 is therefore demonstrated because the system provides a structured
reconstruction path. The baseline pipeline could show that tasks ran and files
existed, but it did not bind the output to the change, release, infrastructure,
code, quality, lineage, and catalogue evidence in one chain. The proposed
architecture does.

\section{Evaluation Summary}
\noindent
The evaluation shows that the proposed architecture satisfies R1--R5 within
the implemented scope. More importantly, each requirement is satisfied in terms
of the auditability definition from Chapter~2. R1 and R4 demonstrate
normativity and accountability by recording the governing change, release,
reference checks, and control outcomes. R2 and R3 demonstrate reliability by
recording artifact hashes, source-control provenance, quality status, and
complete run state. R3 and R5 demonstrate referentiality by linking source,
raw, curated, manifest, metadata, and gold records. R2 and R5 demonstrate
verifiability because the recorded hashes, manifests, metadata references, and
audit-chain fields can be inspected independently.

The result is that auditability is implemented as a set of linked evidence
fields, not as a retrospective explanation. The proposed architecture therefore
meets the thesis objective: it improves ETL governance by making the authority,
integrity, runtime state, control decisions, and output reconstruction path
part of the delivery and execution workflow.
